项目开发初期，小组主要通过腾讯会议、即时通信和共享文档开展协作，但没有在每次会议结束后立即形成统一格式的书面纪要。报告整理阶段，小组依据课程时间节点、腾讯会议历史、即时通信内容、共享文档、代码提交与合并记录以及组员回忆，对九次线上会议进行了补充整理。因此，以下内容是对实际协作过程的事后整理，并非会议录音的逐字转写。仓库提交日期用于核对开发产物，不直接等同于组会日期；腾讯会议历史截图也仅用于说明线上会议工具的使用情况，不单独证明每次会议的全部议题和出席情况。

协作接口主要包括：前端与后端确认请求和错误反馈，后端与数据库核对事务与历史快照，AI 与商品数据模块核对推荐依据，测试与地图模块验证定位失败和多角色履约。上述依赖来自项目实现和分工；具体会议内容以小组补充整理并共同确认的记录为依据。

\subsection{前八次线上会议}
\subsubsection{第一次线上会议（9月1日）：分组、仓库与协作方式}
\begin{enumerate}
  \item 根据课程安排确定项目小组由范文丽、李承文、赵紫怡和仲禹潼四名成员组成，并确认项目仓库地址及成员访问权限。
  \item 确定腾讯会议、微信群、飞书共享文档和 GitHub 的用途；普通文档保存在仓库 \code{docs} 目录，外部共享文档地址记录在根目录 \code{README.md} 中。
  \item 确定基础责任范围：范文丽主责前端与看板，李承文主责后端与红包积分，赵紫怡主责数据库与 AI，仲禹潼主责测试与高德地图。
  \item 约定基础分工不等于独立开发，需求核对、接口联调、缺陷修复和报告整理由全组共同参与。
\end{enumerate}

\subsubsection{第二次线上会议（9月3日）：需求确定与需求文档}
\begin{enumerate}
  \item 结合课程初始工程和课程要求，讨论顾客端、商家端、骑手端和管理端的建设范围。
  \item 梳理用户管理、商家管理、店铺管理、商品管理、购物车、订单管理、配送和后台管理等核心模块。
  \item 明确《软件需求规格说明书》的主体结构，包括项目背景、业务概念、功能性需求、非功能性需求、系统角色、用例模型和主要业务流程。
  \item 要求需求描述、业务流程、UML 图和后续代码实现相互对应；新增设想只有在明确业务规则和影响范围后才进入开发任务。
\end{enumerate}

\subsubsection{第三次线上会议（9月6日）：需求复核与系统建模}
\begin{enumerate}
  \item 汇报需求文档的整理进度，集中检查不同成员撰写内容之间的术语、角色和功能边界。
  \item 讨论系统总体架构图和 UML 总用例图，核对四类角色与各功能模块之间的关系。
  \item 梳理用户、店铺、商品、订单、地址、收藏、评价、钱包和积分等核心业务对象及其关联。
  \item 确认阶段一优先完成基础业务功能，个性化、可视化和智能化内容在基础流程稳定后继续完善。
  \item 约定阶段性任务使用独立分支，功能完成并自测后通过 Pull Request 合并。
\end{enumerate}

\subsubsection{第四次线上会议（9月8日）：中期检查与任务调整}
\begin{enumerate}
  \item 按课程中期检查要求，检查用户管理、商家管理、店铺管理、商品管理和订单管理等阶段一功能。
  \item 对尚未完成、前后端不一致或暂时无法正常运行的功能建立任务清单，并按问题所属模块重新分配工作。
  \item 范文丽继续完善四端页面、导航结构和数据看板；李承文继续完善订单业务与资产功能；赵紫怡继续完善数据库结构和 AI 调用；仲禹潼继续整理测试检查项并推进地图功能。
  \item 明确中期检查后进入功能补全和集中联调阶段；需求变化需要同时检查页面、接口、数据结构和测试用例。
\end{enumerate}

\subsubsection{第五次线上会议（9月10日）：分支合并与功能补充}
\begin{enumerate}
  \item 汇报中期检查后各模块的修改进度，通过 Pull Request 合并已经完成的阶段性功能。
  \item 集中处理代码冲突、依赖关系和配置差异，核对前端请求、后端接口、数据库字段和返回数据结构。
  \item 讨论四端整合方式，要求顾客端、商家端、骑手端和管理端使用统一后端服务并遵守角色权限。
  \item 确认继续完善用户偏好、钱包积分、红包、数据看板、AI 辅助交互和地图服务等扩展内容。
  \item 要求各成员在合并前完成基本自测，并在 Pull Request 中说明主要修改和验证情况。
\end{enumerate}

\subsubsection{第六次线上会议（9月12日）：第一次集中联调}
\begin{enumerate}
  \item 开展集中前后端联调，检查登录、店铺浏览、商品查看、购物车、订单提交、订单处理和后台管理等流程。
  \item 范文丽重点处理页面布局、导航和数据展示；李承文重点处理订单及资产业务；赵紫怡重点处理数据库关联和 AI 服务调用；仲禹潼重点执行功能、边界及地图测试。
  \item 建立统一缺陷登记表，记录问题类型、所属模块、问题描述、复现步骤、截图、修改建议和处理状态。
  \item 约定缺陷修复后重新执行相关业务链路，避免仅验证发生错误的单个页面或接口。
\end{enumerate}

\subsubsection{第七次线上会议（9月15日）：开发完成与测试版本}
\begin{enumerate}
  \item 按课程节点检查全部开发任务，将工作重点由功能扩展转向整合、测试和交付。
  \item 合并成员分支，处理冲突、接口差异和配置问题，检查四端入口、角色权限与页面跳转。
  \item 检查用户操作、店铺与商品管理、购物车、订单生成、订单状态变化和后台管理等主要业务链路。
  \item 检查钱包积分、数据看板、AI 辅助交互和地图服务等特色内容是否具备测试条件。
  \item 生成用于交叉验收的可部署测试版本，并核对《软件需求规格说明书》与当前实现是否一致。
  \item 原则上不再扩大功能范围，后续集中处理交叉验收、缺陷修复和展示准备。
\end{enumerate}

\subsubsection{第八次线上会议（9月16日）：交叉验收总结与课程验收准备}
\begin{enumerate}
  \item 汇总本组在9月16日验收其他小组成果时形成的检查记录，依据对方需求文档核对功能完成情况、业务流程和运行表现。
  \item 将交叉验收意见按需求符合性、功能表现、交互问题和测试证据进行整理，不把其他小组的问题登记为本组项目缺陷。
  \item 继续检查本组项目，准备9月18日课程验收所需的账号、店铺、商品、订单和统计数据，并检查展示环境和备用方案。
  \item 范文丽检查页面、导航和看板；李承文检查后端业务及钱包积分；赵紫怡检查数据库、AI 与异常降级；仲禹潼组织复测并检查地图功能。
  \item 讨论最终展示顺序并进行模拟展示，确认9月18日全组参加课程验收和问题回答。
\end{enumerate}

\subsection{第九次线上会议（9月19日）：验收复盘与报告安排}
\begin{enumerate}
  \item 复盘9月18日课程验收，整理展示过程中被指出的问题和需要进一步完善的内容。
  \item 将业务问题、页面细节和运行稳定性问题补充到缺陷登记表，按成员责任范围安排修复和复测。
  \item 修改完成后重新执行核心业务流程，检查修复内容是否影响其他模块。
  \item 讨论实验报告结构，确定需要整理需求分析、系统设计、项目实现、测试、成员分工、开发过程、AI 使用情况和问题反思。
  \item 收集系统架构图、UML 图、Pull Request、贡献统计、会议记录、共享文档、缺陷登记和四端运行截图等材料。
  \item 明确报告事实应以代码仓库、现有文档和可核对截图为依据，讨论过但没有实现的设想不写成最终成果。
  \item 安排9月27日前完成报告撰写、交叉检查、LaTeX 编译和最终提交。
\end{enumerate}
